\section{Conclusion}

The repetition code showed Alice and Bob that redundancy can protect a message, but it also exposed the central cost: stronger protection can consume more channel uses and reduce the communication rate. Shannon's theorem answers the question created by this trade-off. For the BSC, families of codes can communicate reliably at any rate below $1-h_2(p)$, but not above it. The typical-noise argument explains this threshold by comparing the number of likely outputs surrounding each codeword with the total output space. Mutual information extends the same idea beyond the BSC by measuring how much a channel output reveals about its input.

Shannon's theorem identifies the limit without selecting a particular code. Polar codes provide one explicit response to this construction problem. Their recursive XOR transform reorganizes identical channel uses into synthetic channels with different reliabilities. The BEC examples showed this separation first for two bits and then recursively for four. Information bits are placed in the most reliable positions, while the remaining positions are frozen to values known by both the encoder and decoder. This addresses the three goals posed in the introduction: rates can approach the symmetric capacity $I(W)$, block-error probability under successive-cancellation decoding can approach zero, and both encoding and decoding require $O(n\log n)$ operations rather than an exhaustive search through the codebook.

These conclusions are asymptotic rather than guarantees about every finite code. At finite block lengths, some synthetic channels remain only partly polarized, and basic successive-cancellation decoding need not provide the best practical error performance. Furthermore, the symmetric capacity $I(W)$ equals the ordinary capacity $C$ for the symmetric BSC and BEC considered in this paper, but the two need not agree for every binary-input channel. Polar codes are therefore not the unique or universally best realization of Shannon's theorem. Their significance here is more precise: Shannon identifies what reliable communication permits, while polarization supplies one explicit and computationally efficient way to approach that limit.


